% ============================ WORKSHEET 1 ==================================
\worksheetheading{1}{Classical description of the atom}
\data{$e = \SI{1.6e-19}{\coulomb}$\quad $m_e = \SI{9.1094e-31}{\kilo\gram}$\quad
$m_p = \SI{1.6727e-27}{\kilo\gram}$\quad $m_n = \SI{1.6749e-27}{\kilo\gram}$\quad
$1\ \mathrm{amu} = \SI{1.6605e-27}{\kilo\gram}$\quad $N_A = \SI{6.022e23}{\per\mole}$\quad
$c = \SI{3.00e8}{\metre\per\second}$\quad $k = 1/(4\pi\varepsilon_0) = \SI{8.99e9}{\newton\metre\squared\per\coulomb\squared}$.
Calculator allowed.}

\wq{1}{Who discovered what?}
Match each scientist to the date and to the discovery or idea: Democritus, Dalton, Thomson, Millikan,
Rutherford. Discoveries: (a) the nucleus; (b) the elementary charge; (c) indivisible atoms;
(d) the electron; (e) compounds contain fixed ratios of atoms of given masses.
\begin{solution}
Democritus, 5th century BC: (c).\quad Dalton, 1808: (e).\quad Thomson, 1897: (d).\quad
Millikan, 1909: (b).\quad Rutherford, 1911: (a).
\end{solution}

\wq{1}{Counting particles}
Give the number of protons, neutrons and electrons of $\nuc{12}{6}{C}$, $\nuc{14}{6}{C}$,
$\nuc{35}{17}{Cl}$ and $\nuc{56}{26}{Fe}$. Which two are isotopes?
\begin{solution}
\begin{tabular}{@{}lccc@{}}
 & $p = Z$ & $n = A - Z$ & $e = Z$\\
$\nuc{12}{6}{C}$ & 6 & 6 & 6\\
$\nuc{14}{6}{C}$ & 6 & 8 & 6\\
$\nuc{35}{17}{Cl}$ & 17 & 18 & 17\\
$\nuc{56}{26}{Fe}$ & 26 & 30 & 26
\end{tabular}\\
$\nuc{12}{6}{C}$ and $\nuc{14}{6}{C}$ are isotopes: same $Z = 6$, different $A$.
\end{solution}

\wq{1}{Units of length}
Express a typical atomic size, 100~pm, in m, nm and \AA. Express 1~fm in pm. What is the typical
ratio between the size of an atom and the size of its nucleus?
\begin{solution}
$100\ \mathrm{pm} = 100\times10^{-12}\ \mathrm m = 1\times10^{-10}\ \mathrm m = 0.1\ \mathrm{nm} = 1$~\AA.\quad
$1\ \mathrm{fm} = 10^{-15}\ \mathrm m = 10^{-3}\ \mathrm{pm}$.\quad
Nucleus $10^{-3}$ to $10^{-2}$~pm, so the ratio is $100/10^{-2} = 10^4$ to $100/10^{-3} = 10^5$.
\end{solution}

\wq{1}{Cathode rays}
In Thomson's tube the beam bends towards the positively charged plate. What does this show?
Why was it important that the same deflection was obtained with any gas and any metal?
\begin{solution}
Opposite charges attract, so the beam carries a \textbf{negative} charge. Since the particles are the
same whatever the gas or the metal, they are not specific to one element: they are a component of
\emph{every} atom (the electron). Hence the atom is not indivisible.
\end{solution}

\wq{2}{Ions and isoelectronic species}
Give the number of protons, neutrons and electrons of $\nuc{24}{12}{Mg}^{2+}$, $\nuc{27}{13}{Al}^{3+}$,
$\nuc{19}{9}{F}^{-}$ and $\nuc{40}{20}{Ca}^{2+}$. Which species are isoelectronic, and with which noble
gas?
\begin{solution}
\begin{tabular}{@{}lccc@{}}
 & $p$ & $n$ & electrons\\
$\ce{Mg^2+}$ & 12 & 12 & $12-2 = 10$\\
$\ce{Al^3+}$ & 13 & 14 & $13-3 = 10$\\
$\ce{F-}$ & 9 & 10 & $9+1 = 10$\\
$\ce{Ca^2+}$ & 20 & 20 & $20-2 = 18$
\end{tabular}\\
$\ce{Mg^2+}$, $\ce{Al^3+}$ and $\ce{F-}$ are isoelectronic (10 electrons, like neon, $Z = 10$).
$\ce{Ca^2+}$ has 18 electrons, like argon.
\end{solution}

\wq{2}{Reading Rutherford's experiment}
Give the three observations of the gold-foil experiment and, for each one, what it tells us about the
atom. Why could Thomson's ``plum pudding'' model not explain them?
\begin{solution}
(1) Most $\alpha$-particles pass straight through: the atom is mostly empty.
(2) A few are strongly deflected: there is a small region of concentrated positive charge that repels
the positive $\alpha$-particles.
(3) About 1 in $10^4$ bounces back: this region, the nucleus, is tiny but contains almost all the mass
(a light object could not stop a heavy, fast $\alpha$-particle).
In the plum-pudding model the positive charge and the mass are spread over the whole atom; the electric
force on the $\alpha$-particle would stay weak everywhere and no particle would bounce back.
\end{solution}

\wq{2}{How heavy are electrons?}
(a) Compute $m_p/m_e$. (b) What fraction of the mass of a $\iso{12}{C}$ atom (exactly 12~amu) is due to
its electrons? Conclude.
\begin{solution}
(a) $m_p/m_e = 1.6727\times10^{-27}/9.1094\times10^{-31} = 1836$.\\
(b) $6\,m_e / (12\ \mathrm{amu}) = 6\times9.1094\times10^{-31}/(12\times1.6605\times10^{-27})
= 2.74\times10^{-4}$, i.e.\ 0.027\,\%. The mass of an atom is almost entirely in its nucleus.
\end{solution}

\wq{2}{A stadium-sized atom}
Take an atom of diameter 100~pm and a nucleus of diameter $10^{-3}$~pm. If the nucleus were enlarged
to a marble of diameter 1.0~cm, what would be the diameter of the atom? What fraction of the atom's
volume does the nucleus occupy?
\begin{solution}
Ratio of diameters: $100/10^{-3} = 10^5$. Atom: $10^5\times1.0\ \mathrm{cm} = 10^5\ \mathrm{cm} =
1.0\ \mathrm{km}$. Volumes scale as the cube: $(10^{-5})^3 = 10^{-15}$ of the atom's volume.
\end{solution}

\wq{3}{Charge-to-mass ratios}
(a) Compute $e/m$ for the electron and for the proton. (b) An electron and a proton enter the same
electric field with the same speed. Which one is deflected more, and by what factor (the deflection is
proportional to $q/m$)? (c) Why did this make the electron easy to discover with cathode rays?
\begin{solution}
(a) Electron: $1.6\times10^{-19}/9.1094\times10^{-31} = 1.76\times10^{11}\ \mathrm{C\,kg^{-1}}$.
Proton: $1.6\times10^{-19}/1.6727\times10^{-27} = 9.57\times10^{7}\ \mathrm{C\,kg^{-1}}$.\\
(b) Same $|q|$, so the ratio is $m_p/m_e = 1836$: the electron is deflected 1836 times more (in the
opposite direction).\\
(c) The very large $e/m_e$ gives large, easily measured deflections, and showed Thomson that the
particle was far lighter than any atom.
\end{solution}

\wq{3}{From amu to grams per mole}
One atom of gold has a mass of 196.967~amu. (a) Express it in kg. (b) How many atoms are there in
1.00~g of gold? (c) How many moles is that? Deduce the molar mass of gold and state the general rule
linking amu and $\mathrm{g\,mol^{-1}}$.
\begin{solution}
(a) $196.967\times1.6605\times10^{-27} = 3.2706\times10^{-25}\ \mathrm{kg}$ (Table 1-1 gives
$3.2707\times10^{-25}$).\\
(b) $N = 1.00\times10^{-3}\ \mathrm{kg}/3.2707\times10^{-25}\ \mathrm{kg} = 3.06\times10^{21}$ atoms.\\
(c) $n = N/N_A = 3.06\times10^{21}/6.022\times10^{23} = 5.08\times10^{-3}\ \mathrm{mol}$, so
$M = 1.00\ \mathrm g/5.08\times10^{-3}\ \mathrm{mol} = 197\ \mathrm{g\,mol^{-1}}$.
Rule: a mass of $x$~amu per atom corresponds to $x\ \mathrm{g\,mol^{-1}}$, because
$1\ \mathrm{amu}\times N_A = 1.6605\times10^{-27}\times6.022\times10^{23} = 1.000\times10^{-3}\ \mathrm{kg\,mol^{-1}}$.
\end{solution}

\wq{3}{Identify the species}
Identify each species (symbol, $A$, charge). Use a periodic table for the symbols.
(a) 26 protons, 30 neutrons, 23 electrons. (b) 17 protons, 20 neutrons, 18 electrons.
(c) A neutral atom with $A = 39$ and one more neutron than protons.
(d) An ion with 10 electrons, 8 neutrons and charge $2-$.
\begin{solution}
(a) $Z = 26$: iron; $A = 56$; charge $26 - 23 = +3$: $\nuc{56}{26}{Fe}^{3+}$.\\
(b) $Z = 17$: chlorine; $A = 37$; charge $17 - 18 = -1$: $\nuc{37}{17}{Cl}^{-}$.\\
(c) $Z + (Z+1) = 39$ gives $Z = 19$, potassium: $\nuc{39}{19}{K}$ (20 neutrons).\\
(d) Charge $2-$ means $Z = 10 - 2 = 8$: oxygen, $A = 8 + 8 = 16$: $\nuc{16}{8}{O}^{2-}$.
\end{solution}

\wq{3}{Density of nuclear matter}
A $\iso{12}{C}$ nucleus has a mass of about 12~amu and a radius of about 2.7~fm. (a) Compute its
density. (b) What would be the mass of a teaspoon (5~cm$^3$) of nuclear matter? Compare with water
($\SI{1000}{\kilo\gram\per\metre\cubed}$).
\begin{solution}
(a) $m = 12\times1.6605\times10^{-27} = 1.99\times10^{-26}\ \mathrm{kg}$;
$V = \tfrac43\pi r^3 = \tfrac43\pi(2.7\times10^{-15})^3 = 8.2\times10^{-44}\ \mathrm{m^3}$;
$\rho = m/V = 2.4\times10^{17}\ \mathrm{kg\,m^{-3}}$.\\
(b) $5\ \mathrm{cm^3} = 5\times10^{-6}\ \mathrm{m^3}$, so $m = 2.4\times10^{17}\times5\times10^{-6}
= 1.2\times10^{12}\ \mathrm{kg}$ (over a billion tonnes), about $2\times10^{14}$ times denser than water.
Ordinary matter is light because atoms are almost empty.
\end{solution}

\wq{4}{Millikan's oil drops}
Millikan measured the charges of oil drops: $3.2\times10^{-19}$, $4.8\times10^{-19}$,
$8.0\times10^{-19}$, $6.4\times10^{-19}$ and $1.12\times10^{-18}$~C. (a) Show that these values are
consistent with a single elementary charge and find it. (b) How many extra electrons does each drop
carry? (c) Could a drop carry $2.4\times10^{-19}$~C?
\begin{solution}
(a) Divided by the smallest value, $3.2\times10^{-19}$~C, the charges give 1, 1.5, 2.5, 2 and 3.5: not
all whole numbers, so $3.2\times10^{-19}$~C is not the unit. Half of it, $1.6\times10^{-19}$~C, gives
2, 3, 5, 4 and 7: all whole numbers. The largest charge dividing every value is
$e = 1.6\times10^{-19}$~C.\\
(b) 2, 3, 5, 4 and 7 electrons.\\
(c) No: $2.4\times10^{-19}/1.6\times10^{-19} = 1.5$ is not a whole number. Charge is quantised.
\end{solution}

\wq{4}{Where does the missing mass go?}
(a) Using the masses in amu of the proton (1.0073), the neutron (1.0087) and the electron
(1/1823), compute the total mass of the particles that make up a $\iso{12}{C}$ atom. (b) Compare with
the mass of the atom, 12~amu exactly: compute the mass defect $\Delta m$ in amu and in kg. (c) The
energy that holds the nucleus together is $E = \Delta m\,c^2$. Compute it per atom and per mole, and
compare with a typical chemical bond energy (about $400\ \mathrm{kJ\,mol^{-1}}$). Why can chemical
reactions not change one element into another?
\begin{solution}
(a) $6\times1.0073 + 6\times1.0087 + 6/1823 = 6.0438 + 6.0522 + 0.0033 = 12.0993$~amu.\\
(b) $\Delta m = 12.0993 - 12 = 0.0993$~amu $= 0.0993\times1.6605\times10^{-27} = 1.65\times10^{-28}$~kg.\\
(c) $E = 1.65\times10^{-28}\times(3.00\times10^8)^2 = 1.48\times10^{-11}$~J per atom (about 93~MeV);
per mole, $\times N_A$: $8.9\times10^{12}\ \mathrm{J\,mol^{-1}} = 8.9\times10^{9}\ \mathrm{kJ\,mol^{-1}}$,
about $2\times10^7$ times a chemical bond. Chemical reactions only rearrange electrons; they involve far
too little energy to break or rebuild nuclei, so $Z$ is conserved in chemistry.
\end{solution}

\wq{4}{How small is the gold nucleus?}
In Rutherford's experiment, an $\alpha$-particle (charge $+2e$) with kinetic energy 5.0~MeV
($1\ \mathrm{MeV} = 1.6\times10^{-13}$~J) heads straight for a gold nucleus ($Z = 79$). It stops and
turns back at the distance $d$ where all its kinetic energy has become electric potential energy
$E_p = k\,q_1q_2/d$. (a) Compute $d$. (b) What does this tell you about the size of the nucleus?
Compare with the size of the gold atom (about 140~pm).
\begin{solution}
(a) $E_k = 5.0\times1.6\times10^{-13} = 8.0\times10^{-13}$~J.
$d = k\,(2e)(79e)/E_k = 8.99\times10^9\times2\times79\times(1.6\times10^{-19})^2/8.0\times10^{-13}
= 4.5\times10^{-14}$~m $= 45$~fm.\\
(b) The $\alpha$-particle is repelled before touching the nucleus, so the gold nucleus has a radius
smaller than 45~fm $= 0.045$~pm (its actual radius is about 7~fm). That is at least 3000 times
smaller than the atom (140~pm), consistent with the $10^{-3}$--$10^{-2}$~pm of the handout.
\end{solution}
